\begin{ouverture}[statistique]{13}{Tableaux croisés et fréquences}{Partie 4 · Statistiques et probabilités}
{Variables qualitatives · Marges · Conditions · Filtres}
\objectif{Lire et compléter des tableaux d’effectifs.}
\objectif{Distinguer fréquences marginale et conditionnelle.}
\objectif{Choisir le bon dénominateur selon la question.}
\prerequis{Proportions et pourcentages.}
\end{ouverture}
\pagestyle{otMain}\essentiel\fontsize{11}{15}\selectfont
\essSection{1}{Deux variables qualitatives}
\begin{definition}Une variable qualitative décrit une catégorie. Elle est nominale si les catégories ne sont pas ordonnées (transport, couleur), ordinale si un ordre a un sens (degré de satisfaction). Un tableau croisé compte les individus pour chaque couple de catégories.\end{definition}
\begin{savoirfaire}[Contrôler un tableau]
Chaque case compte une intersection de deux catégories. La somme d’une ligne donne son total marginal ; celle d’une colonne donne l’autre total marginal. Les deux façons de sommer doivent retrouver le même effectif total.
\end{savoirfaire}
\begin{definition}La fréquence marginale de $A$ est $n_A/N$. La fréquence de $C$ parmi les individus de $A$ est $n_{A\cap C}/n_A$, si $n_A>0$. Elle est conditionnelle : le groupe de référence est restreint à $A$.\end{definition}
\begin{exemple}Si 36 personnes de $A$ sont inscrites à $C$ parmi 60 personnes de $A$, la fréquence est $36/60=60\,\%$. Si $C$ compte 90 personnes au total, la fréquence de $A$ parmi $C$ serait $36/90=40\,\%$.\end{exemple}
\begin{attention}« Parmi les inscrits » et « parmi le groupe A » n’annoncent pas le même dénominateur. Commencer par entourer le groupe de référence.\end{attention}
\clearpage
\essSection{2}{Filtrer et comparer}
\begin{savoirfaire}[ET, OU, NON]
Le filtre « $A$ ET $C$ » conserve l’intersection. « $A$ OU $C$ » conserve les individus d’au moins une catégorie, sans compter deux fois ceux des deux. « NON $C$ » conserve le complémentaire dans la population étudiée.
\end{savoirfaire}
\begin{propriete}L’effectif de l’union vérifie $n_{A\cup C}=n_A+n_C-n_{A\cap C}$. Une fréquence globale peut être obtenue en additionnant les effectifs des groupes puis en divisant par le total.\end{propriete}
\begin{savoirfaire}[Tableur ou fichier]
Une ligne représente un individu, une colonne une variable. Filtrer, puis compter les individus retenus. Les fonctions de comptage conditionnel ou un tableau croisé dynamique permettent d’organiser les calculs. Vérifier les modalités manquantes et éviter de mélanger des populations de dates différentes.
\end{savoirfaire}
\begin{attention}Un écart de fréquences décrit une association dans les données. Il ne démontre pas une relation de cause à effet. Des groupes de tailles différentes peuvent aussi rendre une comparaison globale trompeuse.\end{attention}
\otIntertitre{À l’oral}
Dire une fréquence en une phrase complète : catégorie comptée, groupe de référence, valeur. Faire ensuite inverser la question à un camarade et expliquer pourquoi la réponse peut changer.
\clearpage\rubrique{Découverte}
\begin{blocLivre}{ex-13-1}
\begin{exercice}[1]{Lire les marges}
Le tableau décrit deux groupes $A,B$ et l’inscription à un club $C$. Donner le total, l’effectif de $A$, celui de $C$ et celui de $A\cap C$.\begin{center}\renewcommand{\arraystretch}{1.35}\begin{tabular}{c|ccc}\toprule
 & $C$ & $\overline C$ & Total\\\midrule
$A$ & 36 & 24 & 60\\
$B$ & 24 & 36 & 60\\
Total & 60 & 60 & 120\\
\bottomrule\end{tabular}\end{center}
\end{exercice}
\end{blocLivre}
\begin{blocLivre}{ex-13-2}
\begin{exercice}[1]{Une fréquence conditionnelle}
À partir du tableau, donner la fréquence de $C$ parmi $A$, puis celle de $C$ parmi $B$.\begin{center}\renewcommand{\arraystretch}{1.35}\begin{tabular}{c|ccc}\toprule
 & $C$ & $\overline C$ & Total\\\midrule
$A$ & 36 & 24 & 60\\
$B$ & 24 & 36 & 60\\
Total & 60 & 60 & 120\\
\bottomrule\end{tabular}\end{center}
\end{exercice}
\end{blocLivre}
\begin{blocLivre}{ex-13-3}
\begin{exercice}[1]{Inverser la référence}
Dans une enquête, 30 cyclistes sont des élèves de seconde. Il y a 50 cyclistes et 120 élèves de seconde. Calculer la fréquence des secondes parmi les cyclistes et celle des cyclistes parmi les secondes.
\end{exercice}
\end{blocLivre}
\begin{blocLivre}{ex-13-4}
\begin{exercice}[1]{Compléter avec des pourcentages}
Une population compte 200 individus : 80 dans $A$ et 120 dans $B$. Dans $A$, 75\,\% sont satisfaits ; dans $B$, 40\,\% le sont. Construire le tableau croisé satisfait / non satisfait.
\end{exercice}
\end{blocLivre}
\begin{blocLivre}{ex-13-5}
\begin{exercice}[1]{Nominal ou ordinal ?}
Classer les variables « moyen de transport », « avis : mauvais, moyen, bon », « couleur des yeux » et « niveau : débutant, intermédiaire, avancé ».
\end{exercice}
\end{blocLivre}

\clearpage\rubrique{Entraînement}
\begin{blocLivre}{ex-13-6}
\begin{exercice}[2]{Marge et condition}
Avec le tableau ci-dessous, calculer la fréquence globale du club, puis la fréquence du groupe $A$ parmi les non-inscrits.\begin{center}\renewcommand{\arraystretch}{1.35}\begin{tabular}{c|ccc}\toprule
 & $C$ & $\overline C$ & Total\\\midrule
$A$ & 36 & 24 & 60\\
$B$ & 24 & 36 & 60\\
Total & 60 & 60 & 120\\
\bottomrule\end{tabular}\end{center}
\end{exercice}
\end{blocLivre}
\begin{blocLivre}{ex-13-7}
\begin{exercice}[2]{ET, OU et complémentaire}
À l’aide du tableau, calculer les effectifs qui vérifient : $A$ et $C$ ; $A$ ou $C$ ; non $C$.\begin{center}\renewcommand{\arraystretch}{1.35}\begin{tabular}{c|ccc}\toprule
 & $C$ & $\overline C$ & Total\\\midrule
$A$ & 36 & 24 & 60\\
$B$ & 24 & 36 & 60\\
Total & 60 & 60 & 120\\
\bottomrule\end{tabular}\end{center}
\label{fig-ex-13-7}
\figureGeo{croise-venn}
\end{exercice}
\end{blocLivre}
\begin{blocLivre}{ex-13-8}
\begin{exercice}[2]{Du fichier au tableau}
Un fichier de huit individus contient les couples : $(A;O)$, $(B;N)$, $(A;N)$, $(A;O)$, $(B;O)$, $(B;N)$, $(A;N)$, $(B;O)$. Dresser le tableau des effectifs, avec les modalités $O,N$.
\end{exercice}
\end{blocLivre}
\begin{blocLivre}{ex-13-9}
\begin{exercice}[2]{Comparer des taux avec des effectifs}
Une formation $A$ réussit à 18 personnes sur 20 ; une formation $B$ à 72 sur 90. Comparer les fréquences de réussite puis calculer la fréquence globale.
\end{exercice}
\end{blocLivre}
\begin{blocLivre}{ex-13-10}
\begin{exercice}[2]{Des données réelles de l’Insee}
Pour la France hors Mayotte, le RP2023 donne 32987021 hommes et 35107259 femmes, dont respectivement 5934049 et 5674817 ont moins de 15 ans. Calculer la part des moins de 15 ans parmi les hommes, parmi les femmes, puis la part des femmes parmi les moins de 15 ans. Arrondir au dixième de pourcentage.\par\smallskip\legende{Source : Insee, dossier France, POP T4, géographie au 01/01/2026, \url{https://www.insee.fr/fr/statistiques/2011101?geo=FRANCE-1}.}
\end{exercice}
\end{blocLivre}

\clearpage\rubrique{Problèmes}
\begin{blocLivre}{ex-13-11}
\begin{exercice}[3]{Un total à retrouver}
Dans le groupe $A$, 24 personnes pratiquent une activité ; elles représentent 60\,\% du groupe. Le groupe $B$ compte 50 personnes, dont 15 pratiquent l’activité. Déterminer l’effectif de $A$ puis la fréquence globale.
\end{exercice}
\end{blocLivre}
\begin{blocLivre}{ex-13-12}
\begin{exercice}[3]{Une annonce impossible}
Une enquête portant sur 35 personnes affirme que « exactement 30\,\% ont répondu oui ». Cette annonce peut-elle correspondre à un effectif entier ? Comment peut-elle devenir cohérente ?
\end{exercice}
\end{blocLivre}
\begin{blocLivre}{ex-13-13}
\begin{exercice}[3]{Des fréquences identiques}
Une catégorie $C$ concerne 40\,\% du groupe $A$ et 40\,\% du groupe $B$. Montrer que sa fréquence globale vaut aussi 40\,\%, quels que soient les effectifs positifs des deux groupes.
\end{exercice}
\end{blocLivre}
\begin{blocLivre}{ex-13-14}
\begin{exercice}[3]{L’effet d’un recrutement}
Un club compte 30 débutants sur 80 membres. On recrute uniquement des débutants. Combien faut-il en recruter au minimum pour que leur part atteigne au moins 50\,\% ?
\end{exercice}
\end{blocLivre}
\begin{blocLivre}{ex-13-15}
\begin{exercice}[3]{Des groupes à pondérer}
Une offre $X$ compte 9 satisfaits sur 10 dans le groupe 1 et 24 sur 40 dans le groupe 2. Une offre $Y$ compte 32 sur 40 dans le groupe 1 et 5 sur 10 dans le groupe 2. Comparer les offres dans chaque groupe puis globalement.
\end{exercice}
\end{blocLivre}

\clearpage
